\documentclass{article}
\usepackage{amsmath,amssymb,amsthm}

\newtheorem{theorem}{Theorem}
\newtheorem{claim}{Claim}[theorem]
\theoremstyle{definition}
\newtheorem{definition}{Definition}

\newcommand{\N}{\mathbb{N}}

\begin{document}

\begin{definition}
A graph $G$ is \emph{bipartite} if its vertex set splits into two classes
$A$ and $B$ such that every edge joins a vertex of $A$ to a vertex of $B$.
\end{definition}

\begin{theorem}\label{thm:bip}
A graph is bipartite if and only if it contains no cycle of odd length.
\end{theorem}

\begin{proof}
One direction is immediate: walking around a cycle alternates between the
two classes, so every cycle has even length. For the converse we may assume
$G$ is connected and fix a vertex $r$.

\begin{claim}\label{cl:dist}
Colouring each vertex by the parity of its distance from $r$ is proper.
\end{claim}

\begin{proof}
Suppose an edge $uv$ joins two vertices of the same parity. The shortest
paths from $r$ to $u$ and to $v$ together with the edge $uv$ contain a
closed walk of odd length, and every closed walk of odd length contains an
odd cycle. This contradicts the hypothesis.
\end{proof}

By Claim~\ref{cl:dist}, the two parity classes witness that $G$ is
bipartite.
\end{proof}

\begin{theorem}
Every tree with $n \ge 2$ vertices has at least two leaves.
\end{theorem}

\begin{proof}
Consider a longest path $P = v_0 v_1 \cdots v_k$ in the tree.
\begin{enumerate}
  \item The endpoint $v_0$ is a leaf.
  \begin{proof}
  A neighbour of $v_0$ outside $P$ would extend $P$, and a neighbour on $P$
  other than $v_1$ would close a cycle.
  \end{proof}
  \item By symmetry $v_k$ is a leaf as well, and $v_0 \neq v_k$ because
        $n \ge 2$ forces $k \ge 1$.
\end{enumerate}
\end{proof}

\begin{theorem}
For all $m, n \in \N$ the following hold.
\begin{enumerate}
  \item[(a)] $m + n = n + m$;
  \item[(b)] $m \cdot n = n \cdot m$.
\end{enumerate}
\end{theorem}

\begin{proof}
We prove both parts by induction.
\begin{itemize}
  \item[(a)] \begin{proof}[Proof of (a)]
  Induct on $n$, using $m + 0 = m = 0 + m$ for the base case.
  \end{proof}
  \item[(b)] \begin{proof}[Proof of (b)]
  Induct on $n$ again, now using part~(a) and distributivity in the step.
  \end{proof}
\end{itemize}
\end{proof}

\end{document}
